``The local degrees of freedom of higher dimensional pure Chern--Simons theories,''
\emph{Physical Review D} \textbf{53}, 593--596 (1996),
\doi{10.1103/PhysRevD.53.593}, arXiv:hep-th/9506187.

\bibitem{BanadosGarayHenneaux1996}
M. Ba\~nados, L. J. Garay, and M. Henneaux,
``The dynamical structure of higher dimensional Chern--Simons theory,''
\emph{Nuclear Physics B} \textbf{476}, 611--635 (1996),
\doi{10.1016/0550-3213(96)00384-7}, arXiv:hep-th/9605159.

\bibitem{MiskovicTroncosoZanelli2005}
O. Mi\v{s}kovi\'c, R. Troncoso, and J. Zanelli,
``Canonical sectors of five-dimensional Chern--Simons theories,''
\emph{Physics Letters B} \textbf{615}, 277--284 (2005),
\doi{10.1016/j.physletb.2005.04.043}, arXiv:hep-th/0504055.

\bibitem{MacDowellMansouri1977}
S. W. MacDowell and F. Mansouri,
``Unified geometric theory of gravity and supergravity,''
\emph{Physical Review Letters} \textbf{38}, 739--742 (1977),
\doi{10.1103/PhysRevLett.38.739}.

\bibitem{StelleWest1980}
K. S. Stelle and P. C. West,
``Spontaneously broken de Sitter symmetry and the gravitational holonomy group,''
\emph{Physical Review D} \textbf{21}, 1466--1488 (1980),
\doi{10.1103/PhysRevD.21.1466}.

\bibitem{HenneauxTeitelboim1992}
M. Henneaux and C. Teitelboim,
\emph{Quantization of Gauge Systems}
(Princeton University Press, 1992).

\bibitem{ReggeTeitelboim1974}
T. Regge and C. Teitelboim,
``Role of surface integrals in the Hamiltonian formulation of general relativity,''
\emph{Annals of Physics} \textbf{88}, 286--318 (1974),
\doi{10.1016/0003-4916(74)90404-7}.

\bibitem{LeeWald1990}
J. Lee and R. M. Wald,
``Local symmetries and constraints,''
\emph{Journal of Mathematical Physics} \textbf{31}, 725--743 (1990),
\doi{10.1063/1.528801}.

\bibitem{IyerWald1994}
V. Iyer and R. M. Wald,
``Some properties of Noether charge and a proposal for dynamical black hole entropy,''
\emph{Physical Review D} \textbf{50}, 846--864 (1994),
\doi{10.1103/PhysRevD.50.846}, arXiv:gr-qc/9403028.

\bibitem{Hosotani2005}
Y. Hosotani,
``Dynamical gauge symmetry breaking by Wilson lines in the electroweak theory,''
\emph{Annals of Physics} \textbf{313}, 170--216 (2004),
arXiv:hep-ph/0504272.

\end{thebibliography}

\end{document}
